\documentclass[12pt,a4paper]{article}
\usepackage[utf8]{inputenc}
\usepackage[T1]{fontenc}
\usepackage[french]{babel}
\usepackage{graphicx}
\usepackage{geometry}
\usepackage[hidelinks]{hyperref}
\usepackage{xcolor}
\usepackage{fancyhdr}
\usepackage{titlesec}
\usepackage{float}
\usepackage{enumitem}
\usepackage{listings}
\usepackage{booktabs}
\usepackage{array}

% Configuration des marges
\geometry{top=2.5cm, bottom=2.5cm, left=2.5cm, right=2.5cm}
\pagestyle{fancy}
\fancyhf{}
\lhead{DevOps \& MLOps}
\rhead{Rapport de TP 2 -- MLOps}
\cfoot{\thepage}

% Format des titres simples et lisibles
\titleformat{\section}{\Large\bfseries\color{blue!70!black}}{\thesection}{1em}{}
\titleformat{\subsection}{\large\bfseries\color{gray!80!black}}{\thesubsection}{1em}{}
\titleformat{\subsubsection}{\normalsize\bfseries\color{black}}{\thesubsubsection}{1em}{}

% Style simple pour le code Python & YAML & Bash
\definecolor{codebg}{HTML}{F8F9FA}
\lstdefinestyle{simplepython}{
  language=Python,
  basicstyle=\ttfamily\small,
  backgroundcolor=\color{codebg},
  frame=single,
  rulecolor=\color{gray!40},
  breaklines=true,
  columns=fullflexible,
  keepspaces=true,
  showstringspaces=false,
  numbers=none,
  xleftmargin=0.8em,
  xrightmargin=0.8em,
  aboveskip=0.5em,
  belowskip=0.5em,
  literate={é}{{\'e}}1 {è}{{\`e}}1 {ê}{{\^e}}1 {à}{{\`a}}1 {ç}{{\c c}}1 {ù}{{\`u}}1 {É}{{\'E}}1 {ô}{{\^o}}1 {î}{{\^i}}1
}
\lstdefinestyle{simplebash}{
  basicstyle=\ttfamily\small,
  backgroundcolor=\color{codebg},
  frame=single,
  rulecolor=\color{gray!40},
  breaklines=true,
  columns=fullflexible,
  keepspaces=true,
  showstringspaces=false,
  numbers=none,
  xleftmargin=0.8em,
  xrightmargin=0.8em,
  aboveskip=0.5em,
  belowskip=0.5em,
  literate={é}{{\'e}}1 {è}{{\`e}}1 {ê}{{\^e}}1 {à}{{\`a}}1 {ç}{{\c c}}1 {ù}{{\`u}}1 {É}{{\'E}}1
}
\lstset{style=simplepython}

% Commande pour insérer une image avec fallback si elle n'existe pas encore
\newcommand{\imageCapture}[3]{%
  \begin{figure}[H]
    \centering
    \IfFileExists{#1}{%
      \includegraphics[width=#2]{#1}%
    }{%
      \fbox{\parbox[c][4cm][c]{#2}{\centering\ttfamily\small [Capture d'écran : #3\\Fichier attendu : #1]}}%
    }
    \caption{#3}
  \end{figure}
}

% Commande pour insérer un logo s'il existe
\newcommand{\logo}[2]{%
  \IfFileExists{#2}{\includegraphics[width=#1]{#2}}{\fbox{\parbox[c][1.5cm][c]{#1}{\centering\scriptsize logo manquant}}}%
}

% ============================================================
% INFORMATIONS DU PROJET ET DES AUTEURS
% ============================================================
\newcommand{\nomEtudiants}{TRABELSI Saifeddine \quad \& \quad TLILI Achref}
\newcommand{\classeEtudiants}{3 IDL 2}
\newcommand{\titreProjet}{TP MLOps : Du jeu de données au modèle surveillé}
\newcommand{\matiere}{DevOps \& MLOps}
\newcommand{\enseignante}{Dr. Mohamed Sahbi Bahroun}
\newcommand{\anneeUniv}{2026 - 2027}

\begin{document}

% ---------------- PAGE DE GARDE ----------------
\begin{titlepage}
\begin{center}

\begin{minipage}[c]{0.25\textwidth}
\begin{flushleft}
\logo{0.9\textwidth}{images/isi.png}
\end{flushleft}
\end{minipage}
\hfill
\begin{minipage}[c]{0.45\textwidth}
\begin{center}
\fontsize{9pt}{11pt}\selectfont
\textbf{République Tunisienne}\\
\textbf{Ministère de l'Enseignement Supérieur}\\
\textbf{et de la Recherche Scientifique}\\
\textbf{Université de Tunis El Manar}\\
\textbf{Institut Supérieur d'Informatique}\\
\end{center}
\end{minipage}
\hfill
\begin{minipage}[c]{0.25\textwidth}
\begin{flushright}
\logo{0.7\textwidth}{images/utm.png}
\end{flushright}
\end{minipage}

\vspace{1.5cm}

{\Large \textbf{RAPPORT DE TRAVAUX PRATIQUES}}\\[0.7cm]
{\large Matière : \textbf{\matiere}}\\[0.3cm]
{\large Spécialité : \textbf{Ingénierie du Développement Logiciel}}\\[0.3cm]
{\large Classe : \textbf{\classeEtudiants}}\\[1cm]

{\large Élaboré par :}\\[0.3cm]
{\Large \textbf{\nomEtudiants}}\\[1cm]

\noindent\rule{\textwidth}{1.5pt}\\[0.35cm]
{\LARGE \textbf{\titreProjet}}\\[0.3cm]
\noindent\rule{\textwidth}{1.5pt}\\[1.5cm]

\begin{center}
\large
\textbf{Enseignant responsable :} \enseignante \\
\end{center}

\vfill
{\large Année Universitaire \anneeUniv}

\end{center}
\end{titlepage}

% ---------------- CONTENU ----------------
\tableofcontents
\newpage

% ============================================================
\section{Présentation Générale du Projet MLOps}
% ============================================================

\subsection{Contexte et Objectifs}
Ce travail pratique applique les principes fondamentaux du MLOps en industrialisant un modèle d'apprentissage automatique de bout en bout. Plutôt que de traiter des exercices isolés, ce projet implémente un cycle de vie unifié autour d'un cas réel : la classification de vins à partir du jeu de données classique \textit{Wine} (178 vins, 13 variables physico-chimiques et 3 classes de cépages).

Le pipeline MLOps repose sur cinq piliers essentiels :
\begin{enumerate}
  \item \textbf{Versionnement des données et reproductibilité (DVC) :} Gérer les jeux de données volumineux sans encombrer Git, avec un pipeline décrit formellement dans \texttt{dvc.yaml}.
  \item \textbf{Tracking des expériences et registre de modèles (MLflow) :} Tracer systématiquement les hyperparamètres, les métriques et les artefacts de 9 entraînements avec validation croisée, puis promouvoir le meilleur modèle avec l'alias \texttt{@champion}.
  \item \textbf{Serving temps-réel et conteneurisation (FastAPI \& Docker) :} Servir le modèle via une API REST sécurisée et validée par Pydantic, empaquetée dans une image Docker non-root.
  \item \textbf{Surveillance statistique de la dérive (Data Drift) :} Détecter de façon proactive l'altération des flux d'inférence en production à l'aide du \textit{Population Stability Index} (PSI) et du test de Kolmogorov-Smirnov (KS).
  \item \textbf{Automatisation continue CI/CD/CT (GitHub Actions) :} Orchestrer l'ensemble avec des portes de qualité strictes bloquant toute régression en production.
\end{enumerate}

\subsection{Arborescence Finale du Projet}
L'organisation du dépôt \texttt{tp-mlops} respecte la structure modulaire demandée :

\begin{lstlisting}[style=simplebash]
tp-mlops/
├── .github/workflows/ct.yml      # Pipeline CI/CD/CT GitHub Actions
├── api/
│   ├── __init__.py
│   ├── main.py                   # API FastAPI (/predict, /health)
│   ├── Dockerfile                # Conteneurisation minimale et non-root
│   └── requirements-api.txt      # Dependances d'inference
├── data/
│   ├── .gitignore                # Exclusion de vins.csv de Git
│   ├── vins.csv                  # Dataset genere par le pipeline DVC
│   ├── production_normale.csv    # Lot sain de production
│   └── production_derive.csv     # Lot derive de production
├── modeles/
│   ├── .gitignore                # Exclusion du modele lourd de Git
│   ├── modele_final.pkl          # Artefact binaire serialise
│   └── run_id.txt                # Identifiant du run selectionne
├── src/
│   ├── preparer.py               # Etape 1 : extraction des donnees
│   ├── valider.py                # Etape 2 : controle qualite automatique
│   ├── entrainer.py              # Etape 3 : entrainement & tracking MLflow
│   ├── enregistrer.py            # Etape 4 : enregistrement registre MLflow
│   ├── simuler_production.py     # Simulation flux de production
│   ├── surveiller.py             # Calcul PSI et KS (detection de derive)
│   ├── superviser.py             # Sonde de disponibilite HTTP
│   └── evaluer.py                # Porte de qualite (non-regression)
├── tests/
│   └── test_api.py               # Tests unitaires et d'integration pytest
├── dvc.yaml & dvc.lock           # Specification et lock du pipeline DVC
├── params.yaml                   # Hyperparametres versionnes
├── metriques.json                # Metriques du modele courant
├── metriques_champion.json       # Metriques de reference du modele en production
└── requirements.txt              # Versions figees de l'environnement complet
\end{lstlisting}

\newpage

% ============================================================
\section{Partie 1 — Données Versionnées et Pipeline Reproductible (DVC)}
% ============================================================

\subsection{Objectifs et Démarche}
Dans un projet d'apprentissage automatique, le code ne représente qu'une fraction du système. Les données et les dépendances doivent être rigoureusement versionnées. Nous utilisons \textbf{DVC (Data Version Control)} pour suivre les données et orchestrer les étapes de préparation, validation et entraînement.

\subsection{Préparation et Contrôle de Qualité}
Le script \texttt{src/preparer.py} extrait le jeu de données depuis \texttt{sklearn.datasets} et sauvegarde 178 lignes et 14 colonnes dans \texttt{data/vins.csv}.

Le script \texttt{src/valider.py} applique un ensemble de contrôles stricts :
\begin{itemize}[leftmargin=*]
  \item Nombre de colonnes exactement égal à 14 (13 mesures chimiques + la cible).
  \item Volume minimal de 150 lignes pour garantir un apprentissage stable.
  \item Taux de valeurs manquantes inférieur à 5\% par caractéristique.
  \item Vérification de cohérence physique (aucune valeur strictement négative).
  \item Vérification de l'ensemble des étiquettes cibles ($\{0, 1, 2\}$).
\end{itemize}

\begin{lstlisting}[style=simplebash]
# Validation reussie sur le jeu de reference :
$ python src/valider.py data/vins.csv
# -> Aucune erreur, code de sortie 0

# Test d'echec sur un lot corrompu (suppression de l'alcool sur 31 lignes) :
$ python src/valider.py data/corrompu.csv
ERREUR QUALITE : alcohol : 17% de valeurs manquantes
# -> Code de sortie 1 (arret immediat du pipeline)
\end{lstlisting}

\subsection{Description du Pipeline dans \texttt{dvc.yaml}}
Le pipeline décrit formellement les dépendances et les sorties de chaque étape :

\begin{lstlisting}[style=simplebash]
stages:
  preparer:
    cmd: python src/preparer.py
    deps: [src/preparer.py]
    outs: [data/vins.csv]
  valider:
    cmd: python src/valider.py data/vins.csv
    deps: [src/valider.py, data/vins.csv]
  entrainer:
    cmd: python src/entrainer.py
    deps: [src/entrainer.py, data/vins.csv]
    params: [entrainement]
    outs: [modeles/modele_final.pkl]
    metrics:
      - metriques.json:
          cache: false
\end{lstlisting}

Lorsqu'on lance la commande \texttt{dvc repro}, DVC exécute les étapes nécessaires. Dès qu'une étape est à jour, elle est ignorée lors des exécutions suivantes :

\imageCapture{images/dvc_repro.png}{0.85\textwidth}{Exécution reproductible du pipeline avec DVC Repro}

\subsection{Réponses aux Questions de la Partie 1}
\begin{itemize}[leftmargin=*]
  \item \textbf{Question 1 — Que contient Git après ce commit ? Où se trouvent réellement les données ?}\\
  \textit{Réponse :} Git ne stocke aucun fichier volumineux. Il ne contient que le fichier pointeur \texttt{vins.csv.dvc} (qui conserve la somme de contrôle MD5, la taille et le chemin), le fichier \texttt{data/.gitignore} et la configuration DVC. Les données réelles sont enregistrées dans le cache DVC local (\texttt{.dvc/cache}) et sur le remote de stockage distant (\texttt{../dvc-stockage}), rangées sous leur hash MD5.

  \item \textbf{Question 2 — Modifier une valeur de \texttt{data/vins.csv}, puis lancer \texttt{dvc status}. Que se passe-t-il ?}\\
  \textit{Réponse :} \texttt{dvc status} signale que l'empreinte MD5 du fichier ne correspond plus au pointeur : \texttt{data/vins.csv: modified}. On peut annuler cette altération et restaurer la version officielle avec \texttt{dvc checkout data/vins.csv.dvc}, ou bien enregistrer une nouvelle version avec \texttt{dvc add data/vins.csv} suivi d'un commit Git.

  \item \textbf{Question 3 — On modifie seulement \texttt{max\_depth} dans \texttt{params.yaml}. Quelles étapes \texttt{dvc repro} relance-t-il ?}\\
  \textit{Réponse :} DVC inspecte le graphe acyclique de dépendances : seule l'étape \texttt{entrainer} dépend de \texttt{params.yaml} (section \texttt{entrainement}). Les étapes \texttt{preparer} et \texttt{valider} sont sautées (\textit{skipped}), ce qui accélère le développement et évite les calculs redondants.
\end{itemize}

\newpage

% ============================================================
\section{Partie 2 — Tracking des Expériences avec MLflow}
% ============================================================

\subsection{Objectifs et Instrumentation}
Le script \texttt{src/entrainer.py} explore une grille de $3 \times 3 = 9$ combinaisons d'hyperparamètres sur un modèle \texttt{RandomForestClassifier} :
\begin{itemize}[leftmargin=*]
  \item \texttt{n\_estimators} $\in \{50, 100, 200\}$
  \item \texttt{max\_depth} $\in \{3, 6, 10\}$
\end{itemize}

Pour chaque run, le script trace dans MLflow :
\begin{itemize}[leftmargin=*]
  \item L'empreinte MD5 du jeu de données utilisé en tag (\texttt{donnees\_md5}).
  \item Les hyperparamètres (\texttt{n\_estimators}, \texttt{max\_depth}, \texttt{seed}).
  \item Les métriques : accuracy d'entraînement (\texttt{acc\_train}), accuracy de test (\texttt{acc\_test}) et surtout la moyenne de validation croisée 5-plis (\texttt{cv\_mean}).
  \item Le modèle sérialisé avec un exemple d'entrée (\texttt{input\_example}) pour en déduire la signature.
\end{itemize}

Le tracking utilise une base SQLite (\texttt{sqlite:///mlflow.db}), prérequis indispensable pour l'activation du registre de modèles.

\subsection{Tableau Comparatif des 9 Runs}

\begin{table}[H]
\centering
\begin{tabular}{lccccc}
\toprule
\textbf{Run Name} & \textbf{n\_estimators} & \textbf{max\_depth} & \textbf{acc\_train} & \textbf{acc\_test} & \textbf{cv\_mean} \\
\midrule
\texttt{rf-50-3}   & \textbf{50}  & \textbf{3}  & \textbf{1.000} & \textbf{1.000} & \textbf{0.986} \\
\texttt{rf-50-6}   & 50  & 6  & 1.000 & 1.000 & 0.986 \\
\texttt{rf-50-10}  & 50  & 10 & 1.000 & 1.000 & 0.986 \\
\texttt{rf-100-3}  & 100 & 3  & 1.000 & 1.000 & 0.986 \\
\texttt{rf-100-6}  & 100 & 6  & 1.000 & 1.000 & 0.986 \\
\texttt{rf-100-10} & 100 & 10 & 1.000 & 1.000 & 0.986 \\
\texttt{rf-200-6}  & 200 & 6  & 1.000 & 1.000 & 0.986 \\
\texttt{rf-200-10} & 200 & 10 & 1.000 & 1.000 & 0.986 \\
\texttt{rf-200-3}  & 200 & 3  & 1.000 & 1.000 & 0.979 \\
\bottomrule
\end{tabular}
\caption{Comparaison des 9 variantes expérimentales dans MLflow.}
\end{table}

\imageCapture{images/mlflow_compare.png}{0.85\textwidth}{Interface de comparaison des runs dans MLflow}

\subsection{Enregistrement et Alias du Modèle Champion}
Le script \texttt{src/enregistrer.py} enregistre le modèle vainqueur dans le registre sous le nom \texttt{classifieur-vins} et lui affecte l'alias officiel \texttt{@champion} :

\begin{lstlisting}[style=simplebash]
$ python src/enregistrer.py
Successfully registered model 'classifieur-vins'.
Created version '1' of model 'classifieur-vins'.
classifieur-vins v1 -> alias @champion
Recharge depuis le registre : PyFuncModel
\end{lstlisting}

\imageCapture{images/mlflow_registry.png}{0.85\textwidth}{Registre de modèles MLflow avec l'alias @champion}

\subsection{Réponses aux Questions de la Partie 2}
\begin{itemize}[leftmargin=*]
  \item \textbf{Question 1 — Pourquoi six runs ont-ils exactement la même accuracy de test (0,972 ou 1,000) ?}\\
  \textit{Réponse :} Le jeu de test ne comporte que 36 vins (20\% de 178). Une seule erreur de prédiction déplace l'accuracy d'un pas discret de $1/36 \approx 0.0278$. Les valeurs discrètes possibles sont limitées, provoquant de nombreuses égalités. La validation croisée (\texttt{cv\_mean} sur 5 plis) permet une évaluation beaucoup plus stable et continue.

  \item \textbf{Question 2 — Quels runs montrent un sur-apprentissage ? Sur quel indice ?}\\
  \textit{Réponse :} Avec \texttt{max\_depth} $\ge 6$, l'accuracy d'entraînement atteint 1.000 (mémorisation totale). Augmenter la profondeur à 10 n'apporte aucun gain en validation croisée par rapport à une profondeur modérée, traduisant une complexité inutile propice à l'overfitting.

  \item \textbf{Question 3 — Quelle configuration retenir, et pourquoi ?}\\
  \textit{Réponse :} La configuration retenue est \texttt{n\_estimators=50, max\_depth=3} (ou \texttt{n\_estimators=100, max\_depth=6}). Selon le principe du rasoir d'Ockham, à performance égale, on choisit toujours le modèle le plus simple : moins d'arbres et moins de profondeur garantissent des prédictions plus rapides et une empreinte mémoire réduite.

  \item \textbf{Question 4 — Comment revenir au modèle précédent si la version 2 pose problème en production ?}\\
  \textit{Réponse :} Il suffit de déplacer l'alias par script ou API :
  \begin{center}
  \texttt{client.set\_registered\_model\_alias("classifieur-vins", "champion", 1)}
  \end{center}
  Tout service chargeant \texttt{models:/classifieur-vins@champion} basculera immédiatement sur la version 1 au rechargement suivant, sans changer une seule ligne de code ni rebâtir d'image Docker.
\end{itemize}

\newpage

% ============================================================
\section{Partie 3 — Servir le Modèle : API REST et Docker}
% ============================================================

\subsection{Architecture de l'API FastAPI (\texttt{api/main.py})}
L'API REST est construite avec FastAPI et respecte les critères industriels :
\begin{itemize}[leftmargin=*]
  \item \textbf{Chargement unique du modèle :} Le fichier \texttt{modele\_final.pkl} est chargé au démarrage de l'application via \texttt{joblib.load}, évitant une latence disque à chaque requête.
  \item \textbf{Validation stricte avec Pydantic :} Le schéma \texttt{Entree} exige exactement 13 valeurs numériques. Toute entrée incomplète est rejetée immédiatement avec un code HTTP 422.
  \item \textbf{Cohérence des colonnes :} Les 13 valeurs sont injectées dans un \texttt{DataFrame} portant les mêmes noms de colonnes qu'à l'entraînement (\texttt{FEATURES}).
  \item \textbf{Réponse enrichie :} La réponse inclut la classe prédite, le niveau de confiance (\texttt{probabilites.max()}) et la version du modèle.
  \item \textbf{Route de santé (\texttt{/health}) :} Permet aux orchestrateurs (Docker, Kubernetes) de tester la disponibilité du conteneur.
\end{itemize}

\subsection{Validation par les Tests Automatisés (\texttt{tests/test\_api.py})}
Quatre tests couvrent les exigences contractuelles : prédiction valide, bonne classe pour le premier vin de référence, rejet des requêtes incomplètes (422) et état du endpoint \texttt{/health}.

\begin{lstlisting}[style=simplebash]
$ pytest tests/ -v
tests/test_api.py::test_prediction_valide PASSED                         [ 25%]
tests/test_api.py::test_vin_connu_bien_classe PASSED                     [ 50%]
tests/test_api.py::test_entree_incomplete_refusee PASSED                 [ 75%]
tests/test_api.py::test_health PASSED                                    [100%]
======================== 4 passed in 1.84s =========================
\end{lstlisting}

\subsection{Conteneurisation Docker et Prédiction Réelle}
Le fichier \texttt{api/Dockerfile} adopte les bonnes pratiques : image de base allégée (\texttt{python:3.11-slim}), mise en cache ordonnée des dépendances, utilisateur non privilégié (\texttt{appuser}) et exposition du port 8000.

\begin{lstlisting}[style=simplebash]
$ docker build -f api/Dockerfile -t api-modele-ml:1.0 .
$ docker run -d --name api-vins -p 8000:8000 api-modele-ml:1.0

# Requete de test HTTP POST :
$ curl.exe -X POST http://localhost:8000/predict -H "Content-Type: application/json" \
  -d "{\"features\": [13.2, 1.78, 2.14, 11.2, 100, 2.65, 2.76, 0.26, 1.28, 4.38, 1.05, 3.4, 1050]}"
\end{lstlisting}

\textbf{Réponse JSON reçue :}
\begin{lstlisting}[style=simplebash]
{"prediction": 0, "confiance": 0.94, "version_modele": "1"}
\end{lstlisting}

\imageCapture{images/swagger_predict.png}{0.85\textwidth}{Interface Swagger de l'API FastAPI et prédiction réussie}

\subsection{Supervision de Disponibilité (\texttt{src/superviser.py})}
Le script interroge \texttt{/health} toutes les 10 secondes. Lors d'un arrêt provoqué du conteneur (\texttt{docker stop api-vins}), une alerte est émise ; dès son redémarrage (\texttt{docker start api-vins}), le rétablissement est notifié :

\begin{lstlisting}[style=simplebash]
17:00:01 OK
17:00:03 OK
# docker stop api-vins
17:00:06 ALERTE : /health a repondu None (echec n 1)
17:00:08 ALERTE : /health a repondu None (echec n 2)
# docker start api-vins
17:00:16 RETABLI apres 2 echec(s)
17:00:16 OK
\end{lstlisting}

\imageCapture{images/supervision_alert.png}{0.85\textwidth}{Trace de la supervision détectant la panne et le rétablissement}

\subsection{Réponses aux Questions de la Partie 3}
\begin{itemize}[leftmargin=*]
  \item \textbf{Question 1 — Pourquoi l'image pèse-t-elle près de 500 à 700 Mo alors que le modèle ne fait que quelques dizaines de Ko ?}\\
  \textit{Réponse :} Le modèle binaire ne pèse que $\sim 70$ Ko. L'essentiel du poids provient du runtime Python, de l'image de base Debian slim et des bibliothèques scientifiques (NumPy, SciPy, Scikit-learn, Pandas) qui intègrent de lourds binaires C/Fortran d'algèbre linéaire (BLAS/LAPACK).

  \item \textbf{Question 2 — Que se passe-t-il si l'image est construite avec scikit-learn 1.6 alors que le modèle a été entraîné avec la 1.5.2 ?}\\
  \textit{Réponse :} Scikit-learn émet un avertissement \texttt{InconsistentVersionWarning}. Selon les évolutions internes de sérialisation, le modèle peut refuser de se charger (\texttt{UnpicklingError}) ou produire des prédictions divergentes. D'où l'absolue nécessité de figer des versions identiques dans \texttt{requirements.txt} et \texttt{requirements-api.txt}.
\end{itemize}

\newpage

% ============================================================
\section{Partie 4 — Surveillance de la Dérive des Données}
% ============================================================

\subsection{Problématique du Data Drift}
Un modèle en production continue de répondre \texttt{200 OK} même lorsque les données réelles changent radicalement. En l'absence d'étiquettes réelles immédiates, nous mesurons la distance entre les distributions de référence et de production.

\subsection{Métriques Utilisées : PSI et Test KS}
Le script \texttt{src/surveiller.py} calcule pour chaque caractéristique :
\begin{enumerate}
  \item \textbf{Population Stability Index (PSI) :} Découpage en 10 classes quantiles sur la référence. Un seuil supérieur à $0.25$ signale une dérive importante.
  \item \textbf{Test de Kolmogorov-Smirnov (KS) :} Test non paramétrique comparant les distributions cumulées. Une $p$-value $< 0.01$ rejette l'hypothèse d'équivalence.
\end{enumerate}

Si une dérive critique est constatée, le script s'arrête avec le code \textbf{2}, conçu pour alerter automatiquement le pipeline de ré-entraînement continu (CT).

\subsection{Résultats Comparatifs}

\begin{table}[H]
\centering
\begin{tabular}{lcccc}
\toprule
\textbf{Caractéristique} & \textbf{PSI} & \textbf{Statistique KS} & \textbf{p-value} & \textbf{Dérive détectée} \\
\midrule
\multicolumn{5}{l}{\textit{Lot avec dérive volontaire (\texttt{production\_derive.csv})}} \\
\midrule
\textbf{alcohol}          & \textbf{2.482} & \textbf{0.476} & \textbf{2.9e-15} & \textbf{True} \\
\textbf{color\_intensity} & \textbf{0.786} & \textbf{0.353} & \textbf{1.8e-08} & \textbf{True} \\
proanthocyanins           & 0.049          & 0.049          & 9.9e-01          & False \\
proline                   & 0.048          & 0.055          & 9.7e-01          & False \\
magnesium                 & 0.037          & 0.045          & 1.0e+00          & False \\
malic\_acid               & 0.032          & 0.036          & 1.0e+00          & False \\
\bottomrule
\end{tabular}
\caption{Rapport de dérive généré par \texttt{src/surveiller.py}.}
\end{table}

\begin{lstlisting}[style=simplebash]
# Sortie de console sur le lot derive :
ALERTE : derive detectee sur ['alcohol', 'color_intensity'] -> declencher le re-entrainement
# -> Code de sortie 2
\end{lstlisting}

\imageCapture{images/drift_report.png}{0.85\textwidth}{Rapport de surveillance de dérive identifiant les deux variables altérées}

\subsection{Réponses aux Questions de la Partie 4}
\begin{itemize}[leftmargin=*]
  \item \textbf{Question 1 — Pourquoi la feature \texttt{proline} n'est-elle pas détectée alors qu'on lui a ajouté du bruit ?}\\
  \textit{Réponse :} La variable \texttt{proline} s'étend de 278 à 1680 (écart-type naturel $\sigma \approx 315$). Le bruit d'écart-type 30 est négligeable devant cette dispersion naturelle ($\approx 1/10$ de $\sigma$). La distribution globale varie à peine (PSI $\approx 0.048$). Une dérive n'est significative que si son amplitude dépasse la variabilité naturelle de la grandeur.

  \item \textbf{Question 2 — Le lot normal n'a pas un PSI nul. Pourquoi ? Pourquoi le seuil est-il à 0,25 et pas à 0 ?}\\
  \textit{Réponse :} Le lot de production est un sous-échantillon aléatoire (120 vins). En raison des fluctuations naturelles d'échantillonnage, les quantiles ne sont jamais mathématiquement identiques au jeu de référence. Un seuil à 0 déclencherait de fausses alertes permanentes. Les seuils standards (0.10 pour la vigilance, 0.25 pour l'action) absorbent ce bruit statistique.

  \item \textbf{Question 3 — Comment vérifier que la dérive dégrade réellement le modèle ?}\\
  \textit{Réponse :} 
  1. Dès que les étiquettes réelles de terrain sont collectées, recalculer la matrice de confusion et l'accuracy sur ce lot pour mesurer la chute réelle de performance.
  2. En production temps-réel sans étiquettes, monitorer l'entropie des probabilités : une baisse de la confiance moyenne signale que le modèle doute.
  3. Vérifier l'importance des variables concernées : \texttt{alcohol} et \texttt{color\_intensity} étant parmi les variables les plus discriminantes, leur dérive impactera directement les prédictions.
\end{itemize}

\newpage

% ============================================================
\section{Partie 5 — Pipeline CI/CD/CT et Porte de Qualité}
% ============================================================

\subsection{La Porte de Qualité (\texttt{src/evaluer.py})}
La porte de qualité empêche la régression du service de prédiction. Elle compare le modèle nouvellement entraîné avec les performances du modèle en production (\texttt{metriques\_champion.json}) :
\begin{itemize}[leftmargin=*]
  \item Le nouveau modèle doit dépasser un seuil absolu : \texttt{cv\_mean} $\ge 0.90$.
  \item Le nouveau modèle ne doit pas régresser par rapport au champion : \texttt{cv\_mean} $\ge \text{champion} - 0.005$.
\end{itemize}

\begin{lstlisting}[style=simplebash]
# Test 1 : Modele optimal accepte :
$ python src/evaluer.py
cv_mean nouveau = 0.986 | champion = 0.972
ACCEPTE : le nouveau modele peut etre deploye
# -> Code de sortie 0

# Test 2 : Regression provoquee (max_depth=1) :
$ python src/evaluer.py
cv_mean nouveau = 0.585 | champion = 0.972
REFUS : 0.585 < seuil absolu 0.9
# -> Code de sortie 1 (echec et interruption du pipeline)
\end{lstlisting}

\subsection{Workflow GitHub Actions (\texttt{.github/workflows/ct.yml})}
Le pipeline automatisé orchestre la boucle complète grâce à 4 déclencheurs :
\begin{itemize}[leftmargin=*]
  \item \texttt{push} : CI/CD classique lorsque le code ou les hyperparamètres changent.
  \item \texttt{schedule} (cron) : Entraînement continu planifié chaque lundi à 3h UTC.
  \item \texttt{workflow\_dispatch} : Déclenchement manuel par un Data Scientist.
  \item \texttt{repository\_dispatch} (type \texttt{derive-detectee}) : Déclenché par le job de monitoring quand \texttt{surveiller.py} sort avec le code 2.
\end{itemize}

\imageCapture{images/github_actions_pipeline.png}{0.85\textwidth}{Exécution réussie du pipeline CI/CD/CT sur GitHub Actions}

\subsection{Réponses aux Questions de la Partie 5}
\begin{itemize}[leftmargin=*]
  \item \textbf{Question 4 — Que se passe-t-il si le nouveau modèle obtient \texttt{cv\_mean = 0,955} ?}\\
  \textit{Réponse :} Dans \texttt{src/evaluer.py}, $0.955 < 0.972 - 0.005 = 0.967$. La condition de non-régression échoue : le script quitte avec une erreur (code $\ne 0$). Le job \texttt{entrainer-evaluer} échoue, le job \texttt{livrer} (qui possède \texttt{needs: entrainer-evaluer}) n'est jamais exécuté, et le modèle champion actuel reste en production.

  \item \textbf{Question 5 — Pourquoi tolérer un écart de 0,005 avec le champion ?}\\
  \textit{Réponse :} Les algorithmes basés sur des tirages aléatoires et les découpages en plis présentent une faible variance stochastique. Deux entraînements sur les mêmes données peuvent osciller de $\pm 0.002$. Sans cette tolérance de $0.5\%$, un modèle équivalent serait refusé à tort pour un simple bruit de mesure.

  \item \textbf{Question 6 — Quelles améliorations apporter pour un vrai projet industriel ?}\\
  \textit{Réponse :}
  \begin{enumerate}
    \item \textbf{Serveur MLflow distant partagé :} Utiliser une instance hébergée avec base PostgreSQL et bucket S3/GCS pour les artefacts, sécurisée par jetons d'accès.
    \item \textbf{Remote DVC cloud :} Héberger les jeux de données sur un bucket S3 ou GCS d'entreprise avec gestion fine des droits IAM.
    \item \textbf{Déploiement progressif (Canary / Blue-Green) :} Déployer sur Kubernetes avec Argo Rollouts en routant d'abord 5\% du trafic, avec retour arrière automatique en cas de pic d'erreurs 5xx ou de dégradation de la latence p99.
    \item \textbf{Orchestration automatisée du monitoring :} Planifier la détection de dérive avec Apache Airflow ou Prefect pour déclencher automatiquement l'API GitHub Actions.
    \item \textbf{Évaluation avancée :} Intégrer des tests d'équité (fairness) et des benchmarks de latence sous charge dans la porte de qualité.
  \end{enumerate}
\end{itemize}

\newpage

% ============================================================
\section{Conclusion et Synthèse}
% ============================================================
Ce TP MLOps a permis de mettre en œuvre l'intégralité du cycle de vie industriel d'un modèle d'apprentissage automatique :
\begin{itemize}[leftmargin=*]
  \item Les données sont versionnées et protégées contre la corruption avec \textbf{DVC}.
  \item Les entraînements sont rigoureusement tracés et les modèles promus par alias avec \textbf{MLflow}.
  \item L'exposition en production est assurée par une API REST typée \textbf{FastAPI} conteneurisée sous \textbf{Docker}.
  \item La santé du système est surveillée en temps réel tant au niveau de la disponibilité HTTP qu'au niveau statistique de la dérive des distributions (\textbf{PSI \& KS}).
  \item La boucle est entièrement fermée et automatisée par un pipeline \textbf{CI/CD/CT GitHub Actions} équipé d'une porte de qualité garantissant la non-régression.
\end{itemize}

Ce projet démontre comment les pratiques DevOps modernes s'étendent naturellement aux problématiques spécifiques du Machine Learning en garantissant reproductibilité, robustesse et automatisation continue.

\end{document}
